A \co@subspan(M, K).at(i, j)@ expression uses tile coordinates. Without an explicit step, its element origin is \co@(i * M, j * K)@. Specifying a step changes the distance between successive tile origins while leaving the selected tile extent unchanged.

Adding \texttt{.step(sM, sK)} overrides the stride: tiles are spaced \texttt{sM} rows and \texttt{sK} columns apart instead of being packed:

\begin{lstlisting}[language=croqtile,caption={Strided tiles: .subspan, .step, and .at}]
matrix.subspan(16, 16).step(32, 32).at(i, j);
\end{lstlisting}

\begin{figure}[htbp]\centering
\includegraphics[width=\linewidth,height=0.56\textheight,keepaspectratio]{f09-fig5_subspan_step_light.png}
\caption{The small diagram contrasts packed [3, 3] tiles with step(5, 5). In the code, [16, 16] tiles with step(32, 32) leave 16-element gaps.}
\end{figure}

Tile \texttt{(0,0)} starts at \texttt{(0,0)}, but tile \texttt{(1,0)} starts at row 32 (not row 16), leaving a 16-row gap. Omitting \texttt{.step} uses a step equal to the tile size --- the packed case.

A step smaller than the tile extent creates overlapping windows; a larger step leaves gaps. For example, a 16-row window advanced by eight rows shares half its rows with the next window. The step controls the origin sequence, so it should not be confused with a stride between elements inside one tile.
